\chapter{Discussion}\label{ch:m-discussion}

This chapter interprets the measurements in Chapter~\ref{ch:m-results} in terms of when GPU-initiated communication benefits an application and which communication patterns support that benefit. It reads the evidence in the order it was collected. It begins with the pattern microbenchmarks, asking which patterns and topologies each stack leads and where those leads change hands. It then moves to the two application microbenchmarks, \texttt{moe} for routed AI-style exchange and \texttt{cg\_step} for the solver's operation sequence, and asks whether their rankings still follow the patterns they are built from. It finally asks whether either level correlates with complete aCG execution, including the monolithic device-NVSHMEM solver. The chapter then reads GPU-initiated scaling across all patterns at once, separates the portability cost of the SYCL path from the cost of communication itself, answers the research questions, and establishes the limits of those answers, with particular attention to completion dependencies, progress mechanisms, and differences between the evaluated implementations.

\section{Cross-Level Synthesis}

One idea organizes the whole campaign: what performs is not a library but a complete configured path, meaning an operation, an algorithm, a transport, a memory placement, and a completion contract taken together. The evidence is arranged as a ladder in which each level holds the level below it fixed and adds one property that a real application has. Table~\ref{tab:m-evidence-ladder} names the four levels, the property each introduces, and what that property did to the ranking established beneath it. Read downward, it is this chapter's argument in short form: every added property moved a ranking at least once, so no level is a safe proxy for the next, and the useful question is never whether a benchmark is representative in general but which property it still holds fixed.

\begin{table}[htbp]
\centering
\caption{The evidence ladder. Each level adds one property to the level above it, and each added property changed a ranking that the level above had established.}
\label{tab:m-evidence-ladder}
\footnotesize
\begin{tabularx}{\textwidth}{lYY}
\toprule
Level & Property it adds & Effect on the ranking above\\
\midrule
Patterns & Payload size and topology & Leads change hands by regime; an intra-node result does not carry across nodes\\
\texttt{moe} & Routing-dependent, variable-count payloads & Reverses the bulk redistribution ranking in two of three routing cases at \texttt{8n4g}\\
\texttt{cg\_step} & Operation composition and result placement & Moves the lead to the path whose scalar completion contract is cheapest\\
aCG & Irregular peers, local work, overlap, convergence & Agreement survives only where primitive and local work per GPU also match\\
\bottomrule
\end{tabularx}
\end{table}

The pattern-by-pattern results identify candidate regimes, while the two application microbenchmarks and aCG show which properties must remain invariant for those regimes to transfer. MPI retains the startup-sensitive ping-pong and isolated-halo region; persistent NVSHMEM is favorable for large batches of regular neighbor exchanges; NCCL-based paths lead bulk collectives; and tuned MPI/UCC avoids a severe small-message all-to-all default. These are complete configured-path results, not intrinsic orderings of communication models.

Two of those regimes are bounded by topology and by payload distribution rather than by library. Across every pattern, an intra-node result carries little information about a multi-node one: the steady halo rate falls by roughly an order of magnitude at the first inter-node topology and is then flat in node count, whereas bulk allreduce improves as nodes gain GPUs because hierarchical staging shrinks the inter-node payload. The MoE benchmark then removes the equal-block assumption while holding operation, topology, and total volume fixed, and the bulk ranking reverses in two of three routing cases: NCCL is fastest at \texttt{8n4g} under uniform routing and 1.55 times slower than CUDA MPI under locality-biased routing, while a hotspot distribution penalizes the two one-sided paths by about 1.8 times.

The composite benchmark preserves an operation sequence but changes several contracts relative to both the primitives and the solver. Its MPI halo serializes two \texttt{MPI\_Sendrecv} calls, its reduced scalars are not consumed, and the archived OSHMPI path terminates with results in host memory. The complete aCG solver then introduces an irregular peer graph, rank-dependent volume, memory-bound local work, convergence, and backend-specific overlap opportunities. The change from an OSHMPI point-estimate lead in the synthetic sequence to lower representative NCCL solver times therefore accompanies changes in scalar primitive, result placement, peer structure, schedule, and completion.

Two cross-level checks do agree, and both are narrow in the same way. The matched MPI and NCCL one-double reduction intervals differ by 7.8--11.1\% on average over the five measured rank counts (Section~\ref{sec:reduction-agreement}), in the one case where operation, datatype, and completion boundary are the same on both sides. The device-initiated step then reproduces the monolithic solver's strong-scaling shape to within 8.0\% and 12.6\% mean absolute deviation at grid side 8192, and misses it by 582\% and 853\% at side 512 (Section~\ref{sec:cg-device-correlation}). The second check is the more instructive one, because it keeps the benchmark, the backend, and the operation sequence identical and varies only the local work per GPU: a communication benchmark can therefore fail to represent an application without any communication difference at all, purely because the ratio of local work to transfer differs. Both are descriptive agreement checks, not fitted or held-out predictions. A benchmark becomes useful for application selection only when its operation, datatype, peer structure, progress model, completion boundary, overlap opportunity, and local work per GPU match the application path.

\begin{resultbox}[Main experimental result]
Pattern-oriented benchmarks identify candidate regimes, but application selection requires the communication schedule and its completion dependencies. A low primitive cost is useful only to the extent that it reduces exposed application time under an equivalent result state.
\end{resultbox}

\section{GPU-Initiated Scaling with NVSHMEM}\label{sec:gpu-initiated-scaling}

Because the research questions ask when GPU-initiated communication pays off, and NVSHMEM is the only evaluated stack that offers it, its behavior is worth reading across the patterns rather than one at a time. Table~\ref{tab:m-nvshmem-scale} places NVSHMEM against the fastest measured stack at one node and at eight, for every pattern.

\begin{table}[htbp]
\centering
\caption{NVSHMEM latency relative to the fastest measured stack for the same pattern and payload, scaling up within a node and scaling out to eight nodes. Lower is better; $1.0$ means NVSHMEM is the fastest path.}
\label{tab:m-nvshmem-scale}
\footnotesize
\begin{tabularx}{\textwidth}{lYrr}
\toprule
Pattern & Payload & Scale-up & Scale-out\\
\midrule
Ping-pong & 4\,B & $2.6\times$ & $5.6\times$\\
Halo, steady & 16\,MiB aggregate & $1.9\times$ & $1.0\times$\\
Allreduce & 16\,MiB & $1.0\times$ & $2.5\times$\\
All-to-all & 256\,KiB per peer & $1.4\times$ & $1.8\times$\\
MoE, uniform & 32\,MiB per rank & $1.5\times$ & $1.7\times$\\
\texttt{cg\_step} & Side 512 & $1.8\times$ & $2.9\times$\\
\bottomrule
\end{tabularx}
\end{table}

The column pair shows one consistent movement: with the halo as the single exception, every pattern places NVSHMEM further behind the leader when nodes are added than when GPUs are added within a node. The effect is largest where the measurement is latency-bound, from $2.6\times$ to $5.6\times$ in ping-pong, and smallest where it is bandwidth-bound, from $1.4\times$ to $1.8\times$ in all-to-all. Its intra-node bulk allreduce is the fastest path measured, and it loses that position entirely across nodes. The scalar case moves the same way: the four-byte allreduce costs 19.9\,$\mu$s on one node, in line with every other stack, and 88.2\,$\mu$s at 32 GPUs, where the others stay between 46 and 49\,$\mu$s.

The exception identifies the regime that does suit it. The steady-state halo is the only pattern whose NVSHMEM implementation is a persistent device-invoked kernel with neighbor signaling, and it is also the only pattern NVSHMEM leads at eight nodes, at 40.6\,GB/s against 35.7\,GB/s for MPI and flat in node count. Its advantage is therefore tied to sustained, batched, neighbor-structured traffic rather than to device invocation in general. The solver results agree: the monolithic device-NVSHMEM variant gains on blocking OSHMPI as process count grows, but does not overtake NCCL.

Three properties of the evaluated configuration are candidate explanations, and the campaign does not separate them. First, the proxy-assisted path documented in Section~\ref{sec:ibgda-check} adds host involvement per request. This overhead can be particularly visible where requests are small and frequent, which is also where the measured gap is largest. Second, the benchmark dispatches NVSHMEM collectives to NCCL, so the intra-node agreement is a property of that configuration and its disappearance across nodes may be a change of selected path rather than of algorithm. Third, the device-invoked variants change batching, signaling, and kernel organization at the same time as invocation. Separating these requires the matched-initiation controls listed in Section~\ref{sec:priority-experiments}; until then, the favorable regime for device-initiated NVSHMEM on this platform is the one the halo identifies, and scaling out is where it is least competitive.

The scale-out column is therefore as much a statement about this platform's transport as about device initiation. IBGDA is unavailable in the evaluated environment, so every inter-node device-issued request is serviced by a CPU proxy (Section~\ref{sec:ibgda-check}), the same limitation previously reported for Leonardo \cite{trotter_cpu-_2025}. The shape of the table is what a per-request proxy hand-off would produce: the penalty is largest in the latency-bound cells, smallest in the bandwidth-bound ones, and absent in the one pattern that batches many neighbor writes behind a single persistent kernel. Published CG results on this class of system attribute the residual scale-out deficit of their GPU-initiated variant to missing device-side computational kernels and suboptimal NVSHMEM configuration, and expect that variant to become competitive at larger GPU counts once those are addressed \cite{trotter_cpu-_2025}. Removing the proxy hand-off would then be expected to narrow the scale-out column, and to narrow it most where it is now widest. How much it narrows is not measurable from this archive, and it is a property of a cluster's adapter, driver, and runtime configuration rather than of the programming model, so it has to be established by repeating the sweep on clusters with verified GPU-posted networking and a matched schedule, the first direction in Section~\ref{sec:priority-experiments}.

Inside a node no proxy stands in the path: device-issued operations resolve to mapped peer memory across NVLink, and this is where the evaluated device-initiated stack comes closest to the leaders. Its intra-node bulk allreduce is the fastest result measured here, and its remaining intra-node gaps fall between $1.4\times$ and $2.6\times$ rather than reaching the scale-out figures of the same patterns. Recent work indicates that this intra-node regime, rather than scale-out, is where device-side communication currently has the most to give on current hardware: low-latency collectives built on NCCL's device API come within approximately 7\% of an estimated hardware lower bound inside a single NVLink domain \cite{shen_every_2026}. The exploratory intra-node NCCL measurements in Appendix~\ref{app:nccl-preliminary} move in the same direction for small allreduce under symmetric memory, although they do not verify which invocation path produced the timings and therefore corroborate the direction rather than the mechanism. Read together with the halo, the regime in which device invocation repays its cost on this platform is bounded by traffic structure and by communication domain at once: sustained, batched, neighbor-structured exchange, and, for latency-bound collectives, the intra-node NVLink domain.

\section{Portability Cost of the SYCL Path}\label{sec:sycl-discussion}

The second axis of the study changes the compute interface rather than the communication stack, and the measurements separate the two more cleanly than a solver-level ratio can. Over the five patterns that time communication alone, expressing the same exchange through SYCL MPI rather than CUDA MPI costs about 1\% at the median (Section~\ref{sec:sycl-portability}). That is the expected result once the payload path, transport, and completion boundary are identical and only host-side submission and the buffer's allocator change. The composite step is the first case with a visible gap, at $1.28\times$ median, and it is also the first case that contains computation; its single-GPU control reaches $1.48\times$ with inter-rank transfer removed entirely. On this platform the portability cost therefore belongs to kernels and host orchestration, not to the communication path, and a portable interface is inexpensive wherever the operation is large enough to hide its submission and visible where it is not.

This constrains how the solver-level ratios may be read. The reported SYCL/CUDA median solver times of 1.17--1.45 (Section~\ref{sec:acg-results}) are of the same order as the communication-free control, so nothing in the present evidence requires a communication explanation for them. The conservative reading is that the solver gap is dominated by the same kernel and orchestration costs the control isolates, and that the communication path contributes little to it.

\paragraph{Pending: verified aCG SYCL campaign.}
\emph{TBD.} The equal-accuracy aCG SYCL comparison is not yet part of the reconciled archive: the CUDA correctness classification requires the reprocessing described in Section~\ref{sec:acg-method}, so the values in Section~\ref{sec:acg-results} are retained as supporting timings rather than a verified comparison. The completed campaign is expected to supply matched-correctness solver times on the same topology axis as the CUDA solver, together with a phase breakdown that separates kernel time from exposed communication. This section will then state whether the communication-free control fully explains the solver-level ratio, or whether an additional communication cost appears once the SYCL solver scales out. Until then the portability half of RQ3 is answered at microbenchmark and composite level only.

\section{Research Questions Revisited}

\paragraph{RQ1: conditions for beneficial GPU invocation.}
The evaluated persistent, kernel-invoked NVSHMEM path is favorable for large, batched regular halo exchanges, while MPI retains the small-message and isolated-exchange advantage. In aCG, the monolithic device-NVSHMEM implementation improves relative to blocking OSHMPI at larger process counts but remains behind NCCL in the retained representative curves. The MoE hotspot case adds a limit in the other direction: when 80\% of tokens target one rank, the NVSHMEM and OSHMPI put-based paths are about 1.8 times slower than the \texttt{alltoallv} implementations at 32 GPUs, because neither has a collective algorithm available to reorganize the incast. Because persistence, batching, API, signaling, and solver organization change together, the campaign does not isolate an initiation-only speedup. The benefit of GPU-posted networking remains an open question for the follow-up experiments in Section~\ref{sec:priority-experiments}.

\paragraph{RQ2: practical SHMEM trade-offs.}
SHMEM-style paths fit persistent or explicitly managed exchanges when symmetric allocation, remote addressing, notification, and safe reuse align with the application. Those responsibilities remain visible either in application code or in an adapter. The oneCCL/NCCL path retains asynchronous execution and approaches direct NCCL at a bulk all-to-all endpoint, while the oneCCL/OSHMPI staging path becomes approximately 18 times slower than direct OSHMPI for the measured 16\,MiB intra-node allreduce. The NVSHMEM oneCCL branch establishes bootstrap, lifecycle, and staging prerequisites only; it provides no public collective or device-initiation performance result.

\paragraph{RQ3: application transfer and CUDA--SYCL retention.}
Matched MPI and NCCL scalar intervals show approximate cross-level agreement, and the device-initiated step reproduces the monolithic solver's scaling shape at the one grid side whose local work per GPU is comparable. The regular halo and the synthetic OSHMPI result do not predict the complete solver ranking, because peer structure, overlap, primitive, and completion differ. Transfer is therefore conditional rather than general, and the conditions are the ones listed in Section~\ref{sec:sycl-discussion} and in the cross-level synthesis. The separate SYCL report gives median solver-time ratios of 1.17--1.45 relative to CUDA at its four reported GPU counts, but unresolved CUDA correctness classification keeps those values as supporting timings rather than a verified equal-accuracy comparison; Section~\ref{sec:sycl-discussion} reads them against the communication-free control and records the campaign that would settle them.

\section{Threats to Validity}

\paragraph{Internal validity.}
The benchmark implementations use equivalent logical operations but not identical algorithms: NCCL all-to-all is grouped point-to-point, MPI uses a collective, and SHMEM may use puts. The persistent NVSHMEM neighbor kernels also differ from their host-initiated baselines in API, batching, signaling, completion, and sampling granularity, so their stack-level differences cannot be attributed to initiation alone. This is appropriate for stack selection but prevents attribution to the programming model alone. The tuned aggregates do not preserve complete environment manifests, the all-to-all tree merges campaigns collected under different defaults, and oneCCL launch placement is less explicit than the \texttt{srun} mapping. UCC's dispatch is observed in the installed 1.4.0 runtime, while internal algorithm names are checked against the public \texttt{v1.4.x} source branch because no exact public \texttt{v1.4.0} tag is available. The OSHMPI device-path A/B changes operand residency, symmetric allocator, and UCC together, so only its phase-specific observations, not a single-cause interpretation of the total, are supported. MPI regime changes in aCG further weaken large-scale MPI comparisons.

\paragraph{Construct validity.}
Latency and logical bandwidth do not capture every application cost. Steady halo amortizes synchronization, while phase instrumentation deliberately introduces extra synchronization; its sum is diagnostic rather than an alternative headline time. The winning OSHMPI \texttt{cg\_step} result stages two scalar operands through host memory and is therefore a configured-stack result rather than a like-for-like device-resident reduction comparison. The aCG ``halo'' timer measures only exposed wait and under-reports OSHMPI's blocking begin-phase cost. Programmability is assessed through required allocations, synchronization, source integration, and semantic constraints; it is not a controlled developer study or a quantitative software-engineering metric.

\paragraph{External validity.}
The single-platform evaluation described in Section~\ref{sec:eval-environment} does not establish performance portability across GPU generations, accelerator vendors, or other network and software configurations. GPU-posted NVSHMEM RMA remains outside the evaluated transport scope (Section~\ref{sec:ibgda-check}). Only two sparse matrices validate the complete solver. Redistribution is validated only at the microbenchmark level: MoE supplies variable-count, routing-dependent evidence, but no complete application exercises the all-to-all family here.

\paragraph{Conclusion validity.}
Ping-pong, halo, allreduce, and MoE cells normally contain five jobs and fifteen trial records, while the analyzed all-to-all and final \texttt{cg\_step} cells contain three jobs and nine records. The MoE campaign is also a later benchmark snapshot than the reconciled inputs and is archived separately for that reason. No hypothesis tests or confidence intervals are used, so small differences are not treated as definitive. The aCG report retains the reference paper's fastest-converged-trial convention for work-normalized comparison, but Figure~\ref{fig:m-acg-backend-median} makes median completed-trial time the representative backend view because a minimum can favor rare good modes. MPI is additionally shown by per-job minimum, median, and maximum. Host NVSHMEM in aCG disables NCCL collective dispatch, so its difference from OSHMPI cannot be generalized to a fully tuned NVSHMEM configuration.

\paragraph{Preliminary NCCL evidence.}
The newer NCCL study, retained in Appendix~\ref{app:nccl-preliminary}, has only rounded three-run medians and qualitative spread information. Missing source revisions, run commands, datatype, correctness logs, and timing boundaries prevent replication from the summary alone. Comparing 2.18.5 plain with 2.31.2 symmetric changes both version and memory handling; comparing the two 2.31.2 columns still combines allocation and registration. Neither comparison isolates device initiation, and the proposed warm-up and GDAKI failure explanations remain unverified. The study identifies a follow-up candidate after the main scalar-completion and halo-representation gaps are addressed.

These limitations motivate the five research directions in Section~\ref{sec:priority-experiments}: validation on IBGDA-enabled clusters, isolation of initiation effects, stronger benchmark--application correspondence, evaluation of newer device-communication interfaces, and a device-initiated path behind a portable collective interface.

\begin{keypoint}[Scope of the claims]
The results establish which complete configured stacks performed best for the measured communication signatures on Leonardo. They do not establish an intrinsic ordering among the MPI, collective, and PGAS programming models.
\end{keypoint}
